\documentclass[conference]{IEEEtran}

\usepackage[utf8]{inputenc}
\usepackage[T1]{fontenc}
\usepackage[brazil]{babel}
\usepackage{booktabs}
\usepackage{graphicx}
\usepackage{url}

\title{Assistente via Telegram para Consultas Operacionais em Centros de Distribuição: Avaliação de uma Classificação Híbrida de Intenções}

\author{
\IEEEauthorblockN{Daniel Domiciano Spindola, Gustavo de Oliveira Fusco e João Gabriel Coitinho Bonin}
\IEEEauthorblockA{Engenharia de Computação, Fundação Hermínio Ometto (FHO)\\
Araras, Brasil\\
RA: 115562, 115220 e 115872}
}

\begin{document}
\maketitle

\begin{abstract}
Centros de distribuição mantêm informações sobre embarques, documentos fiscais, agendamentos, pátio e ocorrências em fontes distintas, o que torna consultas operacionais dependentes de pesquisa manual. Este trabalho implementa e avalia um assistente via Telegram que combina regras determinísticas, extração de identificadores e classificação supervisionada de intenções. O ambiente operacional usa quatro bancos SQLite com dados sintéticos. O corpus de desenvolvimento contém 2.000 mensagens, balanceadas em dez classes; o conjunto de teste Out-of-Distribution (OOD) contém 210 mensagens sintéticas, incluindo variações linguísticas e casos dirigidos a regras. Foram comparados Multinomial Naive Bayes, LinearSVC calibrado e Random Forest calibrado com TF-IDF de unigramas e bigramas. O LinearSVC foi selecionado por F1-macro em validação cruzada agrupada. No OOD, o classificador isolado com limiar de confiança 0,55 obteve acurácia de 0,690 e F1-macro de 0,707. O Engine híbrido obteve acurácia de 0,648 e F1-macro de 0,658. Portanto, neste conjunto sintético, os resultados não demonstram vantagem da combinação de regras e modelo sobre o classificador isolado. O estudo entrega um protótipo funcional e reprodutível, mas sua avaliação não substitui testes com mensagens reais e anonimizadas.
\end{abstract}

\begin{IEEEkeywords}
classificação de intenções, logística, Telegram, TF-IDF, aprendizado supervisionado, dados sintéticos
\end{IEEEkeywords}

\section{Introdução}

Centros de distribuição coordenam atividades de transporte, agendamento, entrada de veículos, carregamento, emissão de documentos fiscais e registro de ocorrências. Quando essas informações ficam distribuídas em sistemas diferentes, uma pergunta simples sobre uma carga pode exigir consultas manuais em várias fontes. Isso aumenta o tempo de resposta e o retrabalho de operadores.

Estudos sobre inteligência artificial aplicada à logística discutem a oportunidade de automatizar tarefas e integrar informação operacional \cite{klumpp}. Neste projeto, o problema é delimitado à identificação da intenção de uma mensagem e à recuperação de informações sintéticas já estruturadas. O sistema não gera decisões logísticas nem substitui as fontes de registro.

O objetivo é desenvolver e avaliar um assistente via Telegram que aceite perguntas em linguagem natural e consulte dados de embarque, nota fiscal, agendamento, pátio, ocorrência, estadia e transportadora. Como objetivos específicos, foram definidos: integrar fontes SQLite por identificadores operacionais; comparar três classificadores supervisionados; combinar classificação com regras e extração de entidades; verificar autorização por perfil; e medir o desempenho em um conjunto OOD sintético.

O restante do artigo está organizado da seguinte forma: a Seção II apresenta trabalhos relacionados; a Seção III descreve os dados, a arquitetura e o protocolo experimental; a Seção IV apresenta e discute os resultados; e a Seção V conclui o trabalho, incluindo suas limitações e próximos passos.

\section{Trabalhos Relacionados}

Klumpp discute aplicações de inteligência artificial e digitalização na logística e na gestão de cadeias de suprimentos \cite{klumpp}. Esse trabalho contextualiza a relevância da automação, mas não avalia o fluxo específico de classificação de intenções e consulta transacional desenvolvido aqui.

Minaee et al. revisam métodos de classificação de texto, incluindo abordagens de aprendizado profundo \cite{minaee}. Tais métodos são relevantes para tarefas linguísticas complexas, mas o presente escopo prioriza uma solução compacta, local e orientada a um conjunto limitado de intenções.

Joulin et al. apresentam técnicas eficientes de classificação de texto baseadas em representações de palavras e modelos lineares \cite{joulin}. Essa linha de trabalho motiva a avaliação de classificadores leves em texto esparso, embora seus resultados não sejam diretamente comparáveis aos deste projeto, que usa dados sintéticos logísticos.

Larson et al. introduzem um conjunto de avaliação para classificação de intenções e detecção de mensagens fora de escopo \cite{larson}. O trabalho destaca que reconhecer somente as classes conhecidas não basta para um assistente confiável. Neste projeto, uma classe explícita ``fora de escopo'' e um limiar de confiança tratam esse problema, mas os resultados obtidos mostram que a detecção ainda é limitada.

Em relação a esses trabalhos, a contribuição deste protótipo está na integração ponta a ponta de regras, identificadores, autorização e consultas operacionais via Telegram, com comparação explícita entre o classificador isolado e o Engine híbrido. A avaliação é exploratória e não pretende demonstrar superioridade geral sobre os métodos da literatura.

\section{Metodologia}

\subsection{Problema e ambiente de dados}

O sistema recebe uma mensagem e um identificador de chat, determina se o usuário está autorizado, identifica uma intenção e os dados necessários, consulta as fontes e retorna uma resposta curta. A unidade de classificação é a mensagem textual; a variável-alvo é uma entre dez intenções: \texttt{ajuda}, \texttt{saudacao}, \texttt{fora\_escopo}, \texttt{consultar\_embarque}, \texttt{consultar\_agendamento}, \texttt{consultar\_nota\_fiscal}, \texttt{consultar\_veiculo\_patio}, \texttt{consultar\_ocorrencia}, \texttt{consultar\_estadia} e \texttt{consultar\_transportadora}.

Para evitar o uso de informações corporativas ou pessoais, os registros operacionais são sintéticos e criados por \texttt{scripts/seed\_db.py}. O projeto usa SQLite, não SQL Server. Três arquivos representam fontes operacionais; um quarto arquivo guarda os usuários e a auditoria. As relações são feitas por identificadores de embarque, placa e transportadora.

\begin{table}[t]
\centering
\caption{Estrutura e volume inicial dos bancos sintéticos}
\label{tab:db}
\scriptsize
\begin{tabular}{lll}
\toprule
Arquivo & Tabelas & Registros gerados \\
\midrule
\texttt{tms.db} & transportadoras & 30 \\
 & embarques & 500 \\
 & notas fiscais & 1.500 \\
 & agendamentos & 500 \\
\texttt{patio.db} & movimentacao & 500 \\
\texttt{ocorrencias.db} & ocorrencias & 500 \\
\texttt{app.db} & usuarios & 3 perfis de demonstração \\
 & auditoria & cresce durante a execução \\
\bottomrule
\end{tabular}
\end{table}

Os três perfis armazenados em \texttt{usuarios} são interno, transportadora e motorista. O usuário interno pode consultar todas as operações; uma transportadora fica restrita à própria frota; e um motorista, à placa cadastrada. Os testes usam bancos temporários para não depender dos arquivos locais. O script de geração apaga e recria os bancos do diretório escolhido; portanto, os bancos locais não devem ser regenerados sem backup quando contiverem cadastros ou auditoria que se queira preservar.

\subsection{Corpus e exemplos OOD}

O corpus de desenvolvimento está em \texttt{data/corpus\_dev.csv} e contém as colunas \texttt{text}, \texttt{label} e \texttt{group}. São 2.000 mensagens, com 200 exemplos por classe. Os textos são únicos e cada mensagem tem um único rótulo. O campo \texttt{group} identifica o grupo de templates de origem e é usado para evitar que variantes do mesmo template sejam distribuídas entre treino e validação. A abordagem de classificação supervisionada segue a formulação geral discutida por Kotsiantis \cite{kotsiantis}.

O teste OOD está em \texttt{data/corpus\_ood.csv}. Ele tem 210 mensagens: 200 frases sintéticas com variações de linguagem e dez casos específicos dirigidos às regras, como identificadores isolados. A Tabela~\ref{tab:counts} informa a distribuição exata. Embora os exemplos abaixo sejam mensagens efetivamente presentes no arquivo de teste, eles foram escritos para o experimento; não são transcrições de conversas reais de motoristas ou operadores.

\begin{table}[t]
\centering
\caption{Quantidade por classe nos corpora}
\label{tab:counts}
\scriptsize
\begin{tabular}{lrr}
\toprule
Classe & Desenvolvimento & OOD \\
\midrule
Ajuda & 200 & 21 \\
Agendamento & 200 & 20 \\
Embarque & 200 & 22 \\
Estadia & 200 & 20 \\
Nota fiscal & 200 & 23 \\
Ocorrência & 200 & 21 \\
Transportadora & 200 & 20 \\
Veículo no pátio & 200 & 21 \\
Fora de escopo & 200 & 20 \\
Saudação & 200 & 22 \\
\midrule
Total & 2.000 & 210 \\
\bottomrule
\end{tabular}
\end{table}

\begin{table*}[t]
\centering
\caption{Exemplos textuais presentes no conjunto OOD}
\label{tab:ood-examples}
\scriptsize
\begin{tabular}{ll}
\toprule
Classe & Frase do corpus \\
\midrule
Embarque & ``cadê a mercadoria que despacharam'' \\
Agendamento & ``qual janela deram pro caminhão'' \\
Nota fiscal & ``o danfe já tá liberado'' \\
Veículo no pátio & ``a carreta já encostou aí'' \\
Ocorrência & ``teve algum pepino com essa carga'' \\
Estadia & ``quanto passou das cinco horas lá'' \\
Transportadora & ``qual firma pegou esse frete'' \\
Saudação & ``e aí meu amigo'' \\
Fora de escopo & ``qual série você recomenda'' \\
Embarque (regra) & ``status do EMB000010 hj'' \\
\bottomrule
\end{tabular}
\end{table*}

\subsection{Pré-processamento e modelos}

O pré-processamento converte o texto para minúsculas, remove acentos com Unidecode e substitui pontuação por espaços, preservando letras e números. Não são aplicados stemming, lematização ou remoção de stopwords. O \texttt{TfidfVectorizer} usa unigramas e bigramas, \texttt{ngram\_range=(1,2)}, para representar palavras e pares de palavras \cite{salton,manning}.

Foram comparados Multinomial Naive Bayes \cite{mccallum}, LinearSVC \cite{cortes,joachims} e Random Forest \cite{breiman}, em pipelines Scikit-Learn \cite{pedregosa}. Os três classificadores foram calibrados com método sigmoide. O GridSearchCV avaliou \texttt{alpha} em \{0,1; 0,5; 1,0\} para Naive Bayes; \texttt{C} em \{0,5; 1,0; 2,0\} para LinearSVC; e número de árvores em \{80; 150\} e profundidade em \{sem limite; 20\} para Random Forest.

\subsection{Arquitetura híbrida}

O adaptador \texttt{python-telegram-bot} trata \texttt{/start}, \texttt{/ajuda} e mensagens de texto. O \texttt{Engine} verifica o chat antes de consultar dados, normaliza palavras-chave, extrai embarque, placa e nota fiscal por expressões regulares e tenta reconhecer transportadoras por correspondência aproximada. Regras e identificadores são usados antes do classificador; saudações não sobrepõem uma intenção operacional identificada.

Quando a mensagem não pode ser resolvida por regras, o modelo retorna a classe mais provável e sua confiança. O limiar foi fixado em 0,55. Abaixo dele, a mensagem é encaminhada ao fluxo de fora de escopo. Se faltar uma entidade necessária, o Engine guarda temporariamente a intenção do mesmo chat por até cinco minutos. Em caso de nomes de transportadoras ambíguos, apresenta opções numeradas e executa a consulta após a escolha.

O Engine aplica autorização às consultas de embarque, agendamento, nota fiscal, pátio, ocorrência, estadia e transportadora. As respostas são formatadas sem Markdown complexo. A auditoria mascara placas, números de notas fiscais e valores tanto na entrada como no resultado armazenado; sua retenção padrão é de 90 dias, configurável. O cálculo demonstrativo de estadia considera cinco horas de tolerância e taxa configurável de R\$ 60,00 por hora excedente. Essa regra é um parâmetro do protótipo, não uma apuração jurídica ou financeira.

\subsection{Protocolo de validação}

A seleção usa F1-macro, que atribui o mesmo peso às classes \cite{sokolova}. O código cria uma partição agrupada treino/validação por \texttt{StratifiedGroupKFold} de cinco folds, com embaralhamento e semente 42; o primeiro fold é reservado para validação. No treino, o GridSearchCV usa uma validação cruzada agrupada de cinco folds, com semente 43, para ajustar hiperparâmetros e selecionar o modelo. O limiar de 0,55 é fixado antes da avaliação OOD. A varredura de limiares é reportada somente na partição de validação. Depois da seleção, o modelo final é ajustado com todo o corpus de desenvolvimento; o OOD é usado apenas para a avaliação final.

No teste OOD são comparados dois fluxos: (1) ML-only, em que o classificador e o limiar produzem a intenção; e (2) híbrido, que executa o mesmo Engine usado pelo CLI e pelo Telegram, com bancos sintéticos temporários. Foram calculadas acurácia, precisão, revocação, F1-macro, taxas de fallback e matrizes de confusão. O protocolo foi executado em Python 3.13.9, com Scikit-Learn 1.9.1 e pandas 2.3.3.

\section{Resultados}

\subsection{Seleção e comparação dos classificadores}

A Tabela~\ref{tab:model-results} apresenta os resultados dos três classificadores. O LinearSVC calibrado foi selecionado por obter o maior F1-macro na validação cruzada agrupada (0,603). A coluna OOD dos modelos apresenta a classificação direta antes da comparação ponta a ponta com o limiar.

\begin{table}[t]
\centering
\caption{Resultados dos classificadores}
\label{tab:model-results}
\scriptsize
\begin{tabular}{lrrr}
\toprule
Modelo & F1-macro CV & Acurácia OOD & F1-macro OOD \\
\midrule
Multinomial Naive Bayes & 0,555 & 0,710 & 0,709 \\
LinearSVC calibrado & \textbf{0,603} & 0,724 & 0,723 \\
Random Forest calibrado & 0,568 & 0,648 & 0,658 \\
\bottomrule
\end{tabular}
\end{table}

\subsection{Comparação ponta a ponta}

Com o limiar de 0,55, o ML-only obteve acurácia 0,690 e F1-macro 0,707. O Engine híbrido obteve acurácia 0,648 e F1-macro 0,658. O híbrido reduziu de 14,7\% para 9,5\% os casos de mensagens dentro do escopo encaminhadas indevidamente ao fallback, mas aumentou de 45,0\% para 50,0\% as mensagens fora de escopo aceitas como suportadas. Portanto, neste conjunto, a redução de alguns fallbacks não compensou os erros introduzidos pelas regras.

Em uma execução local de 210 mensagens, a latência do Engine foi de 7,580 ms em média e 10,716 ms no percentil 95. Como não houve controle de hardware nem repetição em ambientes distintos, esses valores são apenas uma medição exploratória e não uma garantia de latência em produção.

\subsection{Matrizes de confusão}

Nas Tabelas~\ref{tab:cm-ml} e \ref{tab:cm-hybrid}, as linhas representam a classe real e as colunas, a classe prevista. Os códigos são: AJ (ajuda), AG (agendamento), EMB (embarque), EST (estadia), NF (nota fiscal), OCO (ocorrência), TRA (transportadora), PAT (veículo no pátio), OOS (fora de escopo) e SAU (saudação). A diagonal principal contém os acertos.

\begin{table*}[t]
\centering
\caption{Matriz de confusão do ML-only no OOD (limiar 0,55; 210 mensagens)}
\label{tab:cm-ml}
\scriptsize
\setlength{\tabcolsep}{4pt}
\resizebox{\textwidth}{!}{%
\begin{tabular}{lrrrrrrrrrr}
\toprule
Real $\backslash$ Previsto & AJ & AG & EMB & EST & NF & OCO & TRA & PAT & OOS & SAU \\
\midrule
AJ  & 11 & 0 & 5 & 0 & 0 & 1 & 0 & 0 & 4 & 0 \\
AG  & 0 & 9 & 1 & 1 & 0 & 1 & 1 & 3 & 4 & 0 \\
EMB & 0 & 0 & 13 & 0 & 0 & 0 & 0 & 1 & 8 & 0 \\
EST & 0 & 1 & 0 & 15 & 0 & 0 & 0 & 1 & 2 & 1 \\
NF  & 0 & 0 & 6 & 0 & 16 & 0 & 0 & 0 & 1 & 0 \\
OCO & 0 & 0 & 1 & 0 & 0 & 16 & 1 & 0 & 3 & 0 \\
TRA & 0 & 0 & 1 & 0 & 0 & 0 & 17 & 0 & 2 & 0 \\
PAT & 0 & 0 & 1 & 0 & 0 & 0 & 0 & 19 & 1 & 0 \\
OOS & 2 & 1 & 2 & 1 & 0 & 0 & 1 & 0 & 11 & 2 \\
SAU & 1 & 0 & 0 & 0 & 0 & 0 & 0 & 0 & 3 & 18 \\
\bottomrule
\end{tabular}}
\end{table*}

\begin{table*}[t]
\centering
\caption{Matriz de confusão do Engine híbrido no OOD (210 mensagens)}
\label{tab:cm-hybrid}
\scriptsize
\setlength{\tabcolsep}{4pt}
\resizebox{\textwidth}{!}{%
\begin{tabular}{lrrrrrrrrrr}
\toprule
Real $\backslash$ Previsto & AJ & AG & EMB & EST & NF & OCO & TRA & PAT & OOS & SAU \\
\midrule
AJ  & 11 & 0 & 5 & 0 & 0 & 1 & 0 & 0 & 4 & 0 \\
AG  & 0 & 8 & 2 & 1 & 0 & 1 & 0 & 5 & 3 & 0 \\
EMB & 0 & 0 & 17 & 0 & 0 & 0 & 0 & 1 & 4 & 0 \\
EST & 0 & 2 & 0 & 13 & 0 & 0 & 0 & 2 & 2 & 1 \\
NF  & 0 & 0 & 8 & 0 & 15 & 0 & 0 & 0 & 0 & 0 \\
OCO & 0 & 0 & 5 & 0 & 0 & 13 & 1 & 0 & 2 & 0 \\
TRA & 0 & 0 & 7 & 0 & 0 & 0 & 11 & 1 & 1 & 0 \\
PAT & 0 & 1 & 1 & 0 & 0 & 0 & 0 & 19 & 0 & 0 \\
OOS & 3 & 1 & 2 & 1 & 0 & 0 & 1 & 0 & 10 & 2 \\
SAU & 1 & 0 & 0 & 0 & 0 & 0 & 0 & 0 & 2 & 19 \\
\bottomrule
\end{tabular}}
\end{table*}

\subsection{Análise}

O híbrido acertou mais mensagens de embarque do que o ML-only (17 contra 13), mas passou a classificar como embarque sete mensagens que pertenciam à classe transportadora; também confundiu notas fiscais e ocorrências com embarques. A classe fora de escopo teve precisão 0,36, revocação 0,50 e F1-score 0,42 no Engine. Isso significa que metade das 20 mensagens OOD dessa classe não foi reconhecida como fora de escopo. A matriz confirma que a detecção de mensagens não suportadas continua sendo uma limitação.

Os resultados não comprovam que as regras determinísticas melhoram o desempenho global. Elas são importantes para resolver identificadores e consultas diretas, mas sua prioridade pode produzir erros de intenção em frases ambíguas. A comparação deve ser interpretada com cautela: o OOD tem somente 210 exemplos, é sintético e foi criado pelos próprios autores. Frases com gírias demonstram variedade textual, mas não substituem dados coletados de usuários ou de uma plataforma institucional.

\section{Conclusão}

Foi implementado um protótipo que recebe mensagens pelo Telegram, classifica dez intenções, extrai identificadores, consulta quatro arquivos SQLite e aplica autorização por perfil. Também foram implementados contexto temporário para perguntas incompletas e escolha de transportadora ambígua, cálculo demonstrativo de estadia e auditoria com mascaramento e retenção.

Os objetivos de integração e implementação foram atingidos no ambiente simulado. Entretanto, a tese de que o híbrido melhora a classificação não foi confirmada: no OOD avaliado, o Engine híbrido ficou abaixo do ML-only em acurácia e F1-macro. O resultado atende à necessidade de avaliar a hipótese sem forçar uma conclusão positiva.

As principais limitações são a origem inteiramente sintética dos dados, a pequena dimensão do teste OOD, a baixa detecção de mensagens fora de escopo e a ausência de validação com usuários. Trabalhos futuros devem coletar mensagens reais com autorização e anonimização, documentar sua origem, separar os dados por sessão ou usuário, ampliar o teste independente e revisar as regras que causaram confusão entre embarque, nota fiscal e transportadora. O cálculo de estadia também precisa ser validado por responsáveis operacionais antes de qualquer uso real.

\begin{thebibliography}{13}

\bibitem{klumpp}
M. Klumpp et al., ``Artificial intelligence for logistics and supply chain management,'' in \emph{Digitalization in Supply Chain Management and Logistics}, 2021, pp. 1--25.

\bibitem{minaee}
S. Minaee et al., ``Deep learning based text classification: A comprehensive review,'' \emph{ACM Computing Surveys}, vol. 54, no. 3, pp. 1--40, 2021.

\bibitem{joulin}
A. Joulin, E. Grave, P. Bojanowski, and T. Mikolov, ``Bag of tricks for efficient text classification,'' in \emph{Proc. 15th Conference of the European Chapter of the Association for Computational Linguistics}, 2017, pp. 427--431.

\bibitem{larson}
S. Larson et al., ``An evaluation dataset for intent classification and out-of-scope prediction,'' in \emph{Proc. Conference on Empirical Methods in Natural Language Processing and the 9th International Joint Conference on Natural Language Processing}, 2019.

\bibitem{kotsiantis}
S. B. Kotsiantis, ``Supervised machine learning: A review of classification techniques,'' \emph{Informatica}, vol. 31, no. 3, pp. 249--268, 2007.

\bibitem{salton}
G. Salton and C. Buckley, ``Term-weighting approaches in automatic text retrieval,'' \emph{Information Processing \& Management}, vol. 24, no. 5, pp. 513--523, 1988.

\bibitem{manning}
C. D. Manning, P. Raghavan, and H. Schütze, \emph{Introduction to Information Retrieval}. Cambridge, U.K.: Cambridge University Press, 2008.

\bibitem{mccallum}
A. McCallum and K. Nigam, ``A comparison of event models for naive Bayes text classification,'' in \emph{AAAI-98 Workshop on Learning for Text Categorization}, 1998, pp. 41--48.

\bibitem{cortes}
C. Cortes and V. Vapnik, ``Support-vector networks,'' \emph{Machine Learning}, vol. 20, pp. 273--297, 1995.

\bibitem{joachims}
T. Joachims, ``Text categorization with support vector machines: Learning with many relevant features,'' in \emph{Proc. European Conference on Machine Learning}, 1998, pp. 137--142.

\bibitem{breiman}
L. Breiman, ``Random forests,'' \emph{Machine Learning}, vol. 45, no. 1, pp. 5--32, 2001.

\bibitem{pedregosa}
F. Pedregosa et al., ``Scikit-learn: Machine learning in Python,'' \emph{Journal of Machine Learning Research}, vol. 12, pp. 2825--2830, 2011.

\bibitem{sokolova}
M. Sokolova and G. Lapalme, ``A systematic analysis of performance measures for classification tasks,'' \emph{Information Processing \& Management}, vol. 45, no. 4, pp. 427--437, 2009.

\end{thebibliography}

\end{document}
